\documentclass{scrartcl}
\usepackage[T1]{fontenc}
\usepackage[ngerman]{babel}
\usepackage{booktabs}
\usepackage{tabularx}
\usepackage{siunitx}
\usepackage{multirow}
\usepackage[table]{xcolor} % für Farben in Tabellen
\usepackage{diagbox} %die obere linke Tabellenzelle diagonal teilen
\usepackage{amsmath}
\usepackage[hidelinks]{hyperref} % für Verweise aus innerhalb der Zellen
\usepackage{graphicx}
\definecolor{hellblau}{HTML}{DCEAFB}
\definecolor{hellgruen}{RGB}{220,245,220}

\sisetup{locale=DE}
\begin{document}
\section{Messreihe}

%TODO: Zellen in Farben
%TODO: Zelle in 2 Dreiecken teilen
%TODO: ref from table cell to some picture

\begin{table}[h] 
	\centering
	\caption{Table with multi-Columns}
	\begin{tabular}{ccc} %
		
		\input{tabellenInhal.tex}

	\end{tabular} 

\end{table} 


\begin{table}[htbp]
  \centering
  \label{tab:messung}
  \begin{tabularx}{\linewidth}{X S[table-format=2.0] S[table-format=2.1]}
    \toprule
    {Messpunkt} & {Zeit (min)} & {Temperatur (\si{\degreeCelsius})} \\
    \midrule
    Start   & 0  & 65.0 \\
    Mitte   & 10 & 52.4 \\
    Ende    & 20 & 44.1 \\
    \bottomrule
  \end{tabularx}
  \caption{Temperaturabfall während des Versuchs}
\end{table}

\begin{table}[htbp]
	\centering
	\caption{Experimental Results}
	\label{tab:results} %TODO: wie verwendet man das auf der 5. Seite
	\begin{tabular}{l c c}
		\hline
		Sample & Temp & Yield \\
		\hline
		A & 100 & 85 \\
		\hline
	\end{tabular}
\end{table}


%\begin{table}[h!]
\begin{table}[h!]	
	\centering
	\caption{Experimental Results}
	\begin{tabular}{||c c c c||} 
		\hline
		Col1 & Col2 & Col2 & Col3 \\ [0.5ex] 
		\hline\hline
		1 & 6 & 87837 & 787 \\ 
		2 & 7 & 78 & 5415 \\
		3 & 545 & 778 & 7507 \\
		4 & 545 & 18744 & 7560 \\
		5 & 88 & 788 & 6344 \\ [1ex] 
		\hline
	\end{tabular}
\end{table}


%\usepackage{multirow}
\begin{table}
	\centering
	\caption{Tabelle mit multi-Spalten (i.e.multi-Columns)}
	\begin{tabular}{|l|l|l|} %TODO: 'l' vs 's' vs 'c'
	
		\hline
		\multicolumn{2}{|c|}{Multi-Column} & Column 3 \\
		\hline
		\multirow{2}{*}{Multi-Row} & Row 1 & Data 1 \\
		\cline{2-3}
		& Row 2 & Data 2 \\
		\hline
	\end{tabular}   
\end{table}


%\usepackage{booktabs}
\begin{table}
	\centering
	\caption{Table with multi-Columns}

\begin{tabular}{lrr}
	\toprule
	\textbf{Product} & \textbf{Q1} & \textbf{Q2} \\
	\midrule
	Product A & \$125,000 & \$142,000 \\
	Product B & \$98,500 & \$105,000 \\
	\bottomrule
\end{tabular}   
\end{table}   

%TODO: https://nhigham.com/2019/11/19/better-latex-tables-with-booktabs/
%% First version of table.
%\begin{tabular}{|l|c|c|c|c|c|l|}\hline
%	& \multicolumn{3}{c|}{$\tol=\tols$} & \multicolumn{3}{c|}{$\tol=\told$}\\
%	& $mv$  & Rel.~err & Time    & $mv$  & Rel.~err & Time   \\\hline
%	\trigmv    & 11034 & 1.3e-7 & 3.9 & 15846 & 2.7e-11 & 5.6 \\
%	\trigexpmv & 21952 & 1.3e-7 & 6.2 & 31516 & 2.7e-11 & 8.8 \\
%	\trigblock & 15883 & 5.2e-8 & 7.1 & 32023 & 1.1e-11 & 1.4e1\\
%	\expleja   & 11180 & 8.0e-9 & 4.3 & 17348 & 1.5e-11 & 6.6 \\\hline
%\end{tabular}
%
%% Second version of table, with booktabs.
%\begin{tabular}{lcccccl}\toprule
%	& \multicolumn{3}{c}{$\tol=\tols$} & \multicolumn{3}{c}{$\tol=\told$}
%	\\\cmidrule(lr){2-4}\cmidrule(lr){5-7}
%	& $mv$  & Rel.~err & Time    & $mv$  & Rel.~err & Time\\\midrule
%	\trigmv    & 11034 & 1.3e-7 & 3.9 & 15846 & 2.7e-11 & 5.6 \\
%	\trigexpmv & 21952 & 1.3e-7 & 6.2 & 31516 & 2.7e-11 & 8.8 \\
%	\trigblock & 15883 & 5.2e-8 & 7.1 & 32023 & 1.1e-11 & 1.4e1\\
%	\expleja   & 11180 & 8.0e-9 & 4.3 & 17348 & 1.5e-11 & 6.6 \\\bottomrule
%\end{tabular}

\begin{table}
	\centering
	\caption{Bunte Tabelle}  
	\begin{tabular}{|l|c|c|}
		\hline
		\rowcolor{blue!20}
		Messung & Temperatur & Status \\
		\hline
		A & 20 & \cellcolor{green!25}OK \\
		\hline
		B & 45 & \cellcolor{yellow!35}Warnung \\
		\hline
		C & 80 & \cellcolor{red!25}Kritisch \\
		\hline
	\end{tabular}  
\end{table}   

\begin{table}
	\centering
	\caption{Ganze Spalte einfärben} 
  
\begin{tabular}{
		l
		>{\columncolor{blue!10}}c
		r
	}
	Messung & Wert & Einheit \\
	A & 12.5 & m \\
	B & 15.8 & m
\end{tabular}  
\end{table}   


\begin{table}
	\centering
	\caption{Abwechselnd gefärbte Zeilen} 
	\rowcolors{2}{gray!15}{white}

\begin{tabular}{lcc}
	Messung & Wert 1 & Wert 2 \\
	A & 10 & 20 \\
	B & 15 & 25 \\
	C & 30 & 40 \\
	D & 35 & 45
\end{tabular}

\rowcolors{2}{}{}
\end{table}   
  
  \begin{table}
  	\centering
  	\caption{Die obere linke Tabellenzelle diagonal geteilt}

  	\begin{tabular}{|l|c|c|}
  		\hline
  		\diagbox[width=3cm,height=1.4cm]
  		{Produkt}{Monat}
  		& Januar & Februar \\
  		\hline
  		Gas   & 120 & 150 \\
  		\hline
  		Strom & 80  & 95  \\
  		\hline
  		Wärme & 60  & 75  \\
  		\hline
  	\end{tabular}
  \end{table}   
 


\begin{table}[htbp]
	\centering
	
	\begin{tabular}{lll}
		\toprule
		Objekt & Beschreibung & Verweis \\
		\midrule
		Energie & Masse-Energie &
		Gleichung~\eqref{eq:energie} \\
		Grafik & Schematische Darstellung &
		Abbildung~\ref{fig:beispiel} \\
		\bottomrule
	\end{tabular}
	\caption{Verweise aus einzelnen Zellen zum Bild oder Formel}
	\label{tab:referenzen}
\end{table}  
  
\begin{figure}[htbp]
	\centering
	\includegraphics[width=10cm]{ContrainedPendulaState.jpg}
	\caption{Schematische Beispielabbildung}
	\label{fig:beispiel}
\end{figure}


\begin{equation}
	E = mc^2
	\label{eq:energie}
\end{equation}  
  
\end{document}
